\documentclass[11pt,a4paper]{article}
\usepackage[T1]{fontenc}
\usepackage[utf8]{inputenc}
\usepackage{amsmath,amssymb,xcolor}
\usepackage[margin=25mm]{geometry}
\usepackage[colorlinks=true,urlcolor=blue]{hyperref}
\definecolor{nexorablue}{HTML}{2459B5}
\setlength{\parindent}{0pt}
\setlength{\parskip}{7pt}
\setlength{\emergencystretch}{2em}
\renewcommand{\familydefault}{\sfdefault}
\begin{document}

{\small\bfseries\color{nexorablue} NEXORA CPGE}\par
{\small MPSI / NEX-CPL-06 / Énoncé}\par
{\LARGE\bfseries Centre d’une composition de similitudes}\par
\vspace{0.25em}\hrule\vspace{0.6em}

Soient $S(z)=iz+2-i$ et $T(z)=2z-1+i$. Calculer $T\circ S$, puis donner son centre, son rapport et son angle orienté.


\end{document}
